\documentclass[11pt]{article}
\usepackage[margin=.9in]{geometry}
\usepackage{amsmath,amssymb,amsthm,mathtools}
\usepackage[T1]{fontenc}
\usepackage{lmodern,microtype}
\usepackage[colorlinks=true,linkcolor=blue,urlcolor=blue]{hyperref}
\newtheorem{theorem}{Theorem}
\newtheorem{lemma}[theorem]{Lemma}
\newcommand{\E}{\mathbb E}
\newcommand{\norm}[1]{\lVert #1\rVert}
\newcommand{\poly}{\operatorname{poly}}
\title{Quadratic variance occupation for Gaussian Fock formulas}
\author{ChatGPT-6 Astra \and Kaarel H\"anni \and Dmitry Vaintrob}
\date{26 October 2026}
\begin{document}
\maketitle
\begin{abstract}
For equal-diagonal Gaussian tanh feature maps in compatible coordinates,
we construct exact unbiased scalar arithmetic formulas with trace second
moment at most two and expected coordinate size
$CT_0\exp[C\sum_j q_j^2]$. Higher Hermite degrees are created rarely,
with probability independent of the physical variance $q_j$; their small
coefficients supply the quadratic variance charge. Compatible
Hilbert--Schmidt perturbations add their norm sum to the exponent.
For a free near-critical Gaussian chain this gives the feature-representation
scale $\exp[O(1+\delta^2H)]$. This is not a posterior learning theorem,
and no variable-diagonal or non-Gaussian extension is asserted.
\end{abstract}

\section{Feature map and precise scope}

Let $H$ be a real separable Hilbert space with a fixed countable coordinate
basis. For $q>0$, put $F(q)=\E\tanh^2(\sqrt q Z)$, $Z\sim N(0,1)$,
and expand in orthonormal probabilists' Hermite polynomials:
\[
 \tanh(\sqrt q Z)=\sum_{k\ {\rm odd}}a_k(q)H_k(Z),\qquad
 b_k(q)=a_k(q)/\sqrt{F(q)}.
\]
For $\norm u=1$, define the normalized Fock feature
\[
 G_q(u)=\bigoplus_{k\ {\rm odd}}b_k(q)u^{\otimes k}.
\]
Full ordered tensor coordinates are used. Parseval gives
$\norm{G_q(u)}=1$, and the Gaussian Hermite identity gives
\[
 \langle G_q(u),G_q(v)\rangle
 =\sum_kb_k(q)^2\langle u,v\rangle^k
 =\frac{\E\tanh(\sqrt q Z_u)\tanh(\sqrt q Z_v)}{F(q)},
\]
where the standard Gaussian pair has correlation $\langle u,v\rangle$.
Thus the common $q$ is a physical preactivation variance, equal at every
query. It is not a query-dependent scalar oracle.

\begin{theorem}\label{thm:main}
On a finite input set $X$, consider finitely many steps $G_{q_j}$, optionally
interspersed with linear maps $J_j+K_j$. Each $J_j$ is a signed partial
coordinate injection (no duplicated old coordinates), and $K_j$ is
Hilbert--Schmidt in the same compatible bases. Assume all true intermediate
states have norm at most one, and every state entering $G_{q_j}$ has
norm exactly one. Suppose initial coordinates have unbiased random scalar
formulas with pointwise trace second moment at most two and common
expected size $T_0\ge2$. Then exact unbiased coordinate formulas satisfy
\[
 \sum_i\E Y_i(x)^2\le2\quad(x\in X),\qquad
 \sup_i\E\operatorname{size}(Y_i)
 \le CT_0\exp[C(S+Q_2)],
\]
where $S=\sum_j\norm{K_j}_{\rm HS}$ and $Q_2=\sum_jq_j^2$.
The random shapes are independent of the query, and the constants are
universal. In particular $S=0$ for an unmodified Fock chain.
\end{theorem}

Formulas here are finite real arithmetic trees with arbitrary real edge
coefficients. A copied coordinate changes its final edge coefficient;
successive coefficients are combined algebraically. A zero is one leaf.
This prevents a new unary node at every skipped layer. Converting a final
weighted tree into ordinary scalar multiplications costs only a constant
per edge. No finite-bit or computational-complexity conclusion is needed
for the theorem. An arbitrary dense basis rotation is not a free copy.

\section{The exact Hermite estimate and scalar factory}

\begin{lemma}\label{lem:tail}
For every $q>0$, $\sum_{k\ge3,\ k\ {\rm odd}}k^2b_k(q)^2\le6q^2$.
\end{lemma}
\begin{proof}
Gaussian Sobolev Parseval and bounded derivatives give
\[
 \sum_k k(k-1)a_k(q)^2
   =q^2\E\bigl[\tanh''(\sqrt q Z)^2\bigr].
\]
For $k\ge3$, $k^2\le\tfrac32k(k-1)$. Also
$\tanh''z=-2\tanh z\,\operatorname{sech}^2z$, so
$\tanh''(z)^2\le4\tanh^2z$. Division by $F(q)>0$ proves the claim.
It is an exact all-degree estimate, not a Taylor truncation.
\end{proof}

Suppose current coordinate oracles $X_i$ have means $u_i$, $\norm u=1$,
pointwise trace second moment at most $V\ge1$, and common expected size
at most $T\ge2$. At an odd degree $k\ge3$, set
\[
 m_k=\lceil kV\rceil,\qquad p_k=(C_*k^2V)^{-1}.
\]
For an ordered word $I=(i_1,\ldots,i_k)$, return zero with probability
$1-p_k$, and otherwise return
\[
 \frac{b_k(q)}{p_k}\prod_{a=1}^k\overline X_{i_a},
\]
where each bar averages $m_k$ independent child copies. Distinct factors
are independent even if their coordinate indices coincide. The oracle
is unbiased. Unit norm of the true mean gives
\[
 \sum_i\E\overline X_i^2
 \le1+\frac{V-1}{m_k},\qquad
 \left(\sum_i\E\overline X_i^2\right)^k\le e.
\]
There is a universal arithmetic-accounting constant $C_0$ for which the
nonzero branch has expected size at most $C_0km_kT$. Since $m_k\le2kV$,
the whole coordinate costs at most
\[
 1+C_0p_kkm_kT\le1+(2C_0/C_*)T\le T
\]
when $C_*\ge4C_0$. At degree one return the coefficient-weighted copy
$b_1(q)X_i$, also of size at most $T$.

Tonelli's theorem and independence give the degree-$k$ trace bound
$e b_k(q)^2/p_k$. Summing and applying Lemma~\ref{lem:tail} yields
\begin{equation}\label{eq:step}
 V_{\rm new}\le b_1(q)^2V+
 e\sum_{k\ge3}\frac{b_k(q)^2}{p_k}
 \le V(1+C_1q^2),\qquad T_{\rm new}\le T.
\end{equation}
All sums are nonnegative, so countably many coordinates cause no exchange
of conditionally convergent series. The higher-degree coordinates are
new orthogonal coordinates: their zero branch does not evaluate an old
coordinate. This is why $p_k$ can be independent of $q$. Averaging before
forming the tensor product also avoids the potentially divergent naive
series $\sum_k b_k(q)^2V^k$.

\section{Variance resets and compatible controls}

Start by averaging two copies if needed, so the trace bound is at most
$3/2$, at constant size cost. For a small step with $C_1q^2\le1/8$,
update the deterministic bound by \eqref{eq:step}. If it first exceeds
$7/4$, it is still below two. Average two independent copies of each
coordinate. Since the true mean has norm one, the trace becomes at most
$(V_{\rm new}+1)/2\le3/2$, at size multiplier at most three.
Between these resets the trace must increase by at least $1/4$, while
each increase is at most $Cq^2$. Thus there are at most $C Q_2+C$ such
resets. A large step starting from $V\le2$ has
$V_{\rm new}\le2(1+C_1q^2)$. Averaging
$\lceil2V_{\rm new}\rceil$ copies restores trace at most $3/2$.
Its size multiplier is polynomial in $1+q^2$, with logarithm $O(q^2)$
above the fixed small-step cutoff. This proves the theorem when $S=0$.

For completeness, the compatible linear-control estimate has the same
reset form. Let $\kappa=\norm K_{\rm HS}$. In row $i$ put
$c_i=\sum_jK_{ij}^2$ and sample $j$ with probability $K_{ij}^2/c_i$.
The oracle $Z_i=K_{ij}X_j/(K_{ij}^2/c_i)$ has mean $(Ku)_i$ and
\[
 \E Z_i^2\le c_iV,\qquad \sum_i\E Z_i^2\le\kappa^2V.
\]
Zero rows are omitted. For small nonzero $\kappa$, add this row oracle
to the copied $J$ coordinate only on a Bernoulli branch of probability
$p=\kappa$, dividing its coefficient by $p$. Fresh row and copy children
are independent. The mean is $(J+K)u$; its trace is bounded by
\[
 V+2\norm{Ju}\norm{Ku}+\kappa^2V/p\le V+C\kappa
\]
while $V\le2$ and $\norm u\le1$. Common expected size grows by at most
$1+C\kappa$, since the zero branch is a literal coordinate copy.
For zero $\kappa$ no branch is called. Above the fixed small-control
cutoff use $p=1$ and an averaging reset; trace and size costs are
polynomial in $1+\kappa$, with logarithm $O(\kappa)$.
When an intermediate true mean has norm below one, averaging $m$ copies
still gives trace at most $1+(V-1)/m$. Resets can therefore be charged
to the combined increments $C(\kappa+q^2)$. Multiplying the costs proves
Theorem~\ref{thm:main} with exponent $C(S+Q_2)$.

Every realized tree is finite almost surely: there are finitely many
trajectory layers, every selected tensor degree is finite, and every
averaging count is finite. No vector subcomputation is shared.

\section{Readouts, near-critical scale, and limitations}

For $v\in\ell_2$, sample coordinate $i$ with probability
$v_i^2/\norm v^2$ and return $\norm v^2Y_i/v_i$. This is an unbiased
oracle for $\langle v,h(x)\rangle$, with second moment at most
$2\norm v^2$ and the same common expected size bound. The zero readout is
trivial. If the readout has global signed margin one on $X$, average
$O(1+\norm v^2)$ independent copies per group. Chebyshev's inequality
makes a group accurate to $1/2$ with probability at least $3/4$.
Use $O(\log(2|X|))$ independent groups; concentration and a union bound
make a majority correct at every input with positive probability.
Markov's inequality simultaneously bounds total formula size by a
constant times its expectation. Thus a fixed-seed classifier exists as
a majority of signs of real arithmetic formulas of total size
\[
 \poly(\log(2|X|),T_0,1+\norm v)\exp[C(S+Q_2)].
\]
This exact real-constant statement is separate from any finite-bit
completion.

For an uncontrolled equal-diagonal Gaussian tanh chain with
$t=1+\delta$, $0\le\delta\le\delta_0$ for a fixed small $\delta_0$,
and bounded initial variance, the free variance comparison gives
$q_j\le C[\delta+(j+1)^{-1}]$. Indeed $F(q)\le q-cq^2$ on a fixed
bounded variance interval: $(q-F(q))/q^2$ is positive and continuous,
with limit two at zero. Thus
$q_{j+1}=t^2F(q_j)\le(1+C\delta)q_j-cq_j^2$.
For a sufficiently large fixed $K$, the interval $[0,K\delta]$ is
invariant when $\delta_0$ is small. Above this interval the sequence
obeys $q_{j+1}\le q_j-(c/2)q_j^2$, so its reciprocal increases by at
least $c/2$ until it enters the interval. This proves the comparison.
Consequently
\[
 Q_2=\sum_{j<H}q_j^2\le C(1+\delta^2H).
\]
The normalized Gaussian Fock feature/readout representation therefore
has resource $\exp[O(1+\delta^2H)]$, multiplied by the stated initial
and readout resources. This identifies a quadratic free-feature
timescale; it is not a learning statement for a margin-conditioned prior.

The common $q_j$ and compatible bases are substantive hypotheses.
Query-dependent variances require input-dependent Hermite coefficients;
the known variance-oracle and normalization helpers still have an
$O(q_j)$ charge. An arbitrary dense change of basis is not a signed
coordinate copy. Margin conditioning may change Gaussian row laws, and
physical readout normalization may change its norm. Nothing here
provides a variable-diagonal or non-Gaussian posterior bridge, or settles
the original canonical Boolean tanh depth question.
\end{document}
